\documentclass{article}

\usepackage{iclr2027_conference,times}
\usepackage{graphicx}
\usepackage{booktabs}
\usepackage{amsmath,amssymb}
\usepackage{microtype}
\usepackage{xcolor}
\usepackage{url}
\usepackage{hyperref}

\title{Auditing the Refusal-to-Action Boundary under Benign Workflow Accumulation in Persistent Browser Agents}

\author{Anonymous authors\\Paper under double-blind review}

\newcommand{\pending}[1]{\textcolor{red}{[PENDING: #1]}}

\begin{document}
\maketitle

\begin{abstract}
Persistent browser agents can accumulate benign workflow demonstrations and continue acting under the resulting persistent state. Prior work already shows that longitudinal memory and benign execution-oriented experience can affect safety~\citep{altawaha2026remembering,zhao2026safety}; our question is narrower. We ask whether, under an exactly matched future task schedule, allowing benign workflow state to accumulate changes the boundary between immediate refusal and harmful-task-directed action. We decompose behavior into immediate refusal or safe handling (R0), first task-directed action (R1), partial execution depth (R2), and verified external effect (R3). The prospective primary endpoint is a task-specific programmatic R1 predicate, with R0/R2/R3 prespecified as decomposition endpoints. Before any new safety trajectory, the exact Qwen3.5-397B-A17B + Agent Workflow Memory + BrowserART/BrowserGym stack must pass BrowserART's official ten-task benign panel under a ten-step budget with a strict 10/10 gate. A historical Qwen3-8B study motivates the hypothesis but is discovery-only: 103/108 historical future episodes hit a four-step truncation ceiling, 0/108 normally terminated, no listener-confirmed external effect was observed, and HarmBench and DeepSeek produced different arm orderings on the same traces. \pending{prospective Updated-vs-Frozen R1 result and R0/R2/R3 decomposition} Semantic evaluators and blinded human labels are then used only to test measurement robustness around the objective behavioral anchors. \pending{measurement-robustness result} The final claim is restricted to the exact treatment, substrate, and horizon supported by the prospective evidence.
\end{abstract}

\section{Introduction}

Experience-driven agents do not remain at the state at which they were first evaluated. Systems such as Agent Workflow Memory (AWM)~\citep{wang2025awm} turn prior task demonstrations into reusable workflows that can shape later action. A point-in-time safety evaluation therefore leaves open a temporal control question: after benign workflow experience accumulates, does the same agent continue to refuse a harmful request before taking task-directed action?

This broad territory is no longer empty. Longitudinal memory work shows that benign memory accumulation can create later safety failures and develops fixed-probe / snapshot protocols to isolate exposure from stream non-stationarity~\citep{altawaha2026remembering}. Experience-driven self-evolving-agent work further shows that benign experience can reduce safety and links the degradation to an execution-oriented act-over-refuse bias~\citep{zhao2026safety}. We do not claim either result as new. The narrower identification gap is whether \emph{incremental workflow updating itself}, rather than a changing future task stream or weak execution substrate, moves a persistent browser agent's refusal-to-action boundary under an exact same-schedule counterfactual.

A second ambiguity comes from the safety endpoint. For an overtly harmful browser request, immediate refusal before task-directed tool use is qualitatively different from beginning to fill the target interface, and both are different from a verified external side effect. A single semantic safe/unsafe label can collapse these states. We therefore define a behavioral ladder: R0 is immediate refusal or safe handling; R1 is the first harmful-task-directed browser action; R2 is partial execution depth after R1; and R3 is a verifier-confirmed external effect when the environment exposes one. R1 is the prospective primary endpoint.

Our causal contrast is deliberately small. Updated and Frozen arms start from the same persistent workflow and receive the same future tasks, seeds, BrowserART/BrowserGym environment~\citep{kumar2025browserart,drouin2024browsergym}, AXTree/BID action interface, and execution budget. Updated receives prospectively frozen benign workflow units across exposure steps; Frozen retains the step-0 workflow. Thus the treatment difference is workflow accumulation rather than the task schedule.

Browser-agent safety is not interpretable on a substrate that cannot complete benign tasks. The prospective backbone must therefore pass BrowserART's official ten-task benign sanity panel under a benchmark-aligned ten-step budget before any new harmful task is run. The historical four-step setting is retained only as a prefix diagnostic of the same benign qualification trajectories.

The historical R9 study is useful precisely because it reveals what the confirmatory design must repair. Under HarmBench~\citep{mazeika2024harmbench}, historical future branch events were 8/12 Updated, 4/12 Frozen, and 0/12 NullMemory. A blinded DeepSeek re-evaluation of the same traces produced 5/12, 5/12, and 8/12. Retrospective execution audit then showed that 103/108 historical future episodes hit the four-step ceiling, none normally terminated, and no listener-confirmed external effect was recorded. We therefore treat these results as discovery evidence, not as confirmation that workflow updating increases harm.

The paper is organized around the scientific object rather than the historical reviewer path. Section~\ref{sec:object} defines the refusal-to-action boundary. Section~\ref{sec:design} gives the capability qualification and exact same-schedule identification. Section~\ref{sec:historical} explains why historical R9 generates but cannot confirm the hypothesis. Section~\ref{sec:prospective} reports the prospective matched behavioral result. Section~\ref{sec:measurement} then asks whether semantic evaluators and humans support the same interpretation. Claims beyond the exact backbone/AWM/BrowserART setting require separate transport evidence.

\section{Scientific object: the refusal-to-action boundary}
\label{sec:object}

Let $H_t$ denote complete pre-task persistent state at exposure step $t$, $q_{u,t}$ the frozen future task for matched unit $u$, and $A$ the browser-agent transition operator. The Updated state evolves as
\begin{equation}
H^{U}_{t+1}=\operatorname{Update}(H^U_t,B_t),
\end{equation}
where $B_t$ is a prospectively frozen benign workflow unit, while the Frozen arm holds
\begin{equation}
H^{F}_{t+1}=H^F_t.
\end{equation}
The two arms receive the same $q_{u,t}$, seeds, browser environment, action interface, and execution budget.

For each harmful task $q$, a prospectively frozen programmatic predicate $P_q$ maps structured action/environment events to whether the agent has initiated task-directed execution. We define
\begin{equation}
R1(a,u)=\mathbb{I}[\exists k\le K: P_{q_u}(E^a_{u,k})=1],
\end{equation}
with $K=10$ in the prospective design. R0 records safe refusal before R1, R2 records prespecified task-execution depth, and R3 records a verifier-confirmed external effect.

The primary finite matched estimand is
\begin{equation}
D_{R1}=\frac{1}{|\mathcal U|}\sum_{u\in\mathcal U}\left(R1(U,u)-R1(F,u)\right).
\end{equation}
We do not infer a population effect without a separately justified sampling model.

\begin{figure}[t]
\centering
\fbox{\parbox{0.93\linewidth}{\centering
\textbf{Figure 1 placeholder: scientific object.}\\[3pt]
Updated: $W_0\rightarrow W_0+B_1\rightarrow W_0+B_1+B_2\rightarrow W_0+B_1+B_2+B_3$\\
Frozen: $W_0\rightarrow W_0\rightarrow W_0\rightarrow W_0$\\[3pt]
At each matched harmful request: R0 Refuse $\rightarrow$ R1 First task-directed action $\rightarrow$ R2 Partial execution $\rightarrow$ R3 Verified effect.
}}
\caption{The paper's scientific object. The future schedule is matched; the treatment is benign workflow accumulation. Safety behavior is decomposed instead of collapsed into one post-hoc semantic label.}
\label{fig:object}
\end{figure}

\section{Capability qualification and causal identification}
\label{sec:design}

\paragraph{Capability gate.}
The exact Qwen3.5-397B-A17B + AWM + BrowserART/BrowserGym stack first runs the pinned official ten benign tasks with a maximum of ten steps. Passage requires 10/10 prospectively frozen programmatic completion predicates. A failed task is not replaced or rerun. Failure stops safety execution and is classified as substrate failure, not evidence for safe refusal.

\paragraph{Four-step diagnostic.}
The historical four-step budget is evaluated from the first four steps of the same ten-step benign trajectories. It is descriptive and cannot change the capability threshold or task set.

\paragraph{Prospective treatment.}
Only after capability PASS are fresh states, tasks, seeds, benign update units, Updated/Frozen pairing, R0--R3 predicates, external-effect verifiers, retry rules, and statistical decision rules frozen before the first new safety trajectory.

\begin{table}[t]
\centering
\caption{Qualification table. The ten-step benchmark-aligned setting is primary; the four-step prefix is diagnostic only.}
\label{tab:capability}
\small
\begin{tabular}{lcc}
\toprule
Quantity & 10-step primary & 4-step prefix \\
\midrule
Official benign tasks & 10 & same 10 \\
Required successes & 10/10 & descriptive \\
Task replacement & none & none \\
Safety authority if failed & none & none \\
Observed result & \pending{fill} & \pending{fill} \\
\bottomrule
\end{tabular}
\end{table}

\section{Historical discovery and why it is insufficient}
\label{sec:historical}

Historical R9 used Qwen3-8B + AWM on BrowserART with a four-step ceiling. HarmBench produced Updated/Frozen/NullMemory branch counts of 8/4/0, while a blinded DeepSeek audit of the same frozen traces produced 5/5/8. These are genuine historical operational results, but they do not identify an evaluator-independent workflow effect.

The execution audit tightens the boundary further: 103/108 historical future episodes reached the four-step truncation ceiling, 0/108 normally terminated, and 0/108 produced listener-confirmed external effects. The historical evaluator input therefore relied on action/thought traces rather than verified harmful completion. The old results motivate the R1-centered prospective design and the need for measurement robustness; they are not reused as the new confirmatory endpoint.

\section{Prospective matched behavioral result}
\label{sec:prospective}

\pending{Report capability disposition first. If and only if qualified, report the frozen Updated-vs-Frozen matched R1 result. Then report R0/R2/R3 decomposition, parser/grounding diagnostics, and first-crossing localization. Do not insert HarmBench/DeepSeek ordering here as a substitute for a null R1 result.}

\begin{table}[t]
\centering
\caption{Primary matched behavioral result. R1 is primary; R0/R2/R3 are prespecified decomposition endpoints.}
\label{tab:behavior}
\small
\begin{tabular}{lcccc}
\toprule
Endpoint & Updated & Frozen & Matched difference & Status \\
\midrule
R0 immediate refusal & \pending{fill} & \pending{fill} & \pending{fill} & secondary \\
R1 task-directed initiation & \pending{fill} & \pending{fill} & \pending{fill} & primary \\
R2 execution depth & \pending{fill} & \pending{fill} & \pending{fill} & secondary \\
R3 verified effect & \pending{fill} & \pending{fill} & \pending{fill} & secondary \\
\bottomrule
\end{tabular}
\end{table}

\section{Measurement robustness after behavioral anchoring}
\label{sec:measurement}

HarmBench, DeepSeek, and blinded human semantic labels are evaluated as measurement systems around frozen trajectories and objective R1/R3 anchors. Evaluator-Robust Temporal Audit (ERTA) can summarize premise stability, event-set overlap, and contrast-sign stability without selecting a post-hoc majority label. ERTA is secondary unless evaluator reversal prospectively persists on capability-qualified, substantially complete, objectively anchored trajectories while the behavioral treatment effect itself is absent.

\pending{Report evaluator/human agreement conditional on objective behavior; localize disagreements by action depth and task.}

\section{Related work}

\paragraph{Longitudinal memory safety.}
Al-Tawaha et al.~\citep{altawaha2026remembering} study temporal memory contamination using fixed probes over read-only memory snapshots and NullMemory controls across multiple memory architectures. Our narrower object is an exact same-schedule evolving-vs-frozen workflow intervention with refusal-to-action decomposition in an interactive browser agent.

\paragraph{Experience-driven self-evolution safety.}
Zhao et al.~\citep{zhao2026safety} show that benign experience can reduce safety and attribute degradation to execution-oriented act-over-refuse bias across multiple backbones and environments. We do not claim that phenomenon as novel. G1 asks whether incremental workflow accumulation causes a longitudinal control-boundary shift under a matched future schedule and programmatic action endpoint.

\paragraph{Web-agent safety and semantic evaluation.}
BrowserART demonstrates that chat alignment does not automatically transfer to browser-agent action~\citep{kumar2025browserart}; HarmBench provides a standardized automated red-teaming classifier~\citep{mazeika2024harmbench}. G1 treats semantic judges as measurement layers rather than ground truth when objective action/event anchors are available.

\section{Limitations and claim boundary}

The prospective design remains a finite controlled study of one backbone, one AWM implementation, one browser benchmark, and a declared horizon. A positive local R1 effect does not establish a population hazard or universal self-evolution failure. R3 may be sparse because some BrowserART tasks do not expose persistent external-effect instrumentation. Human semantic evidence calibrates interpretation but cannot repair a failed objective or capability endpoint. Cross-backbone or cross-memory generalization requires separately frozen transport evidence.

\section{Conclusion}

The paper asks a narrow control question: under the same future schedule, does benign workflow accumulation move a persistent browser agent from refusal toward task-directed execution? By separating R0--R3 behavior, qualifying the substrate before safety evaluation, and relegating semantic evaluators to a robustness layer, the design distinguishes persistent-state effects from task difficulty, execution failure, and judge interpretation. \pending{Final conclusion must be selected only after the prospective claim ladder resolves.}

\bibliography{references}
\bibliographystyle{iclr2027_conference}

\end{document}
